\section{Анализ существующих решений и постановка проблемы}
\label{sec:Chapter3} \index{Chapter3}

В предыдущем разделе были рассмотрены ключевые принципы, на которых основано построение современных сканеров уязвимостей. В этом разделе проводится обзор конкретных open-source сканеров, ставших де-факто стандартом отрасли~--- Trivy, Clair и Grype. Несмотря на то, что все сканеры в той или иной степени реализуют принципы, описанные в предыдущем разделе, у каждого из них есть свои особенности, влияющие на процесс сканирования.

\subsection{Trivy}

\textit{Trivy}~\cite{trivy}~--- открытый сканер уязвимостей, разработанный Aqua Security. На сегодняшний день Trivy является одним из наиболее популярных open-source решений в области сканирования уязвимостей в OCI-образах~--- в первую очередь благодаря простоте использования, обширному покрытию источников данных, активной поддержке и интеграциям. Trivy используется в качестве сканера по умолчанию в Harbor, поддерживается крупными CI/CD-платформами и регулярно фигурирует в обзорах безопасности контейнерной инфраструктуры.

Основной функционал Trivy сосредоточен в виде CLI-утилиты, написанной на языке Go. Помимо собственно сканирования OCI-образов, Trivy умеет анализировать файловые системы, репозитории git, артефакты Kubernetes и другие источники, однако наиболее распространенным сценарием его использования является сканирование контейнерных образов.

\subsection{Clair}

\textit{Clair}~\cite{clair}~--- открытый сканер уязвимостей, разработанный командой Quay~--- проекта container registry, входящего в состав Red Hat. Исторически Clair был первым массово используемым open-source сканером и до сих пор остается важным игроком в индустрии. Clair используется в качестве сканера по умолчанию в реестре Quay.io и поддерживается во многих managed-registry. В отличие от Trivy, Clair изначально проектировался не как CLI-утилита, а как сервис: его архитектура ориентирована на работу в режиме демона, к которому обращаются клиенты по HTTP API.

Характерной особенностью Clair является явное разделение индексирования и сопоставления на два независимых сервиса~--- \textit{indexer} и \textit{matcher}. Indexer принимает на вход манифест образа, проводит индексирование и формирует \textit{IndexReport}~--- список обнаруженных компонентов с метаданными. Этот результат сохраняется в собственной базе данных Clair (PostgreSQL) и в дальнейшем переиспользуется: при повторном обращении к ранее проиндексированному образу индексирование заново не производится. Matcher по запросу с дайджестом ранее сформированного IndexReport обращается к базе уязвимостей и формирует итоговый \textit{VulnerabilityReport}. Такая декомпозиция позволяет Clair эффективно поддерживать сценарий, при котором появление новых записей в базе уязвимостей не требует повторного скачивания и анализа содержимого образов~--- достаточно повторного вызова matcher.

\subsection{Grype}

\textit{Grype}~\cite{grype}~--- открытый сканер уязвимостей, разработанный компанией Anchore. Характерная особенность Grype заключается в том, что сам он не выполняет индексирование, а ожидает на входе уже готовый SBOM~--- список компонентов образа. Индексирование при этом возложено на сопутствующий инструмент \textit{Syft}~\cite{syft}, также разрабатываемый Anchore. Полный цикл сканирования образа в этой связке выглядит так: Syft проводит индексирование образа и генерирует SBOM в одном из поддерживаемых форматов (SPDX, CycloneDX, либо собственный формат Syft), после чего Grype принимает на вход полученный SBOM и сопоставляет его с базой уязвимостей.

Такое разделение функциональности на два отдельных инструмента позволяет встраивать индексирование и сопоставление в произвольные цепочки поставки ПО независимо друг от друга. В частности, SBOM может быть сгенерирован разработчиком на этапе сборки и сохранен вместе с образом, а Grype в дальнейшем может многократно применяться к этому SBOM по мере появления новых записей в базах уязвимостей, не обращаясь к самому образу.

\subsection{Доступ к данным слоев и взаимодействие с container registry}

Несмотря на различия в архитектуре и подходах, все рассмотренные сканеры объединены общей моделью работы с container registry:

\begin{enumerate}
\item сканер получает манифест образа из registry;
\item каждый слой, для которого результат анализа не предствавлен в локальном кеше, сканер скачивает \textit{целиком};
\item полученный gzip-поток полностью разжимается, после чего для полученного tar-архива производится индексирование содержащихся в нем файлов;
\end{enumerate}

Так, полная загрузка слоя здесь не является особенностью реализации какого-либо из рассмотренных сканеров~--- она диктуется самим форматом слоя согласно спецификации OCI и объясняется тем, что gzip-сжатие не допускает извлечения отдельных файлов без распаковки до соответствующей позиции.

\subsubsection{Объем избыточно скачиваемых данных}

Чтобы оценить масштаб этой неэффективности, можно проследить следующий простой пример. Типичный образ Linux-приложения весом 200--500 МБ обычно состоит из нескольких слоев следующего вида:

\begin{itemize}
\item базовый Linux-слой (Ubuntu, Debian, Alpine, Fedora и т.п.)~--- 30--100 МБ;
\item слой с системными пакетами зависимостей~--- ещё 50--200 МБ;
\item слой с языковыми пакетами (pip, npm, go и т.п.)~--- 50--300 МБ;
\item слой с приложением~--- от единиц МБ до сотен МБ в зависимости от типа приложения.
\end{itemize}

Как было показано в предыдущем разделе, содержимое, представляющее для сканера непосредственный интерес,~--- файлы баз пакетных менеджеров и манифесты языковых зависимостей~--- в сумме редко превышает единицы мегабайт даже для образов с сотнями пакетов. То есть полезная часть содержимого, в общем случае, составляет лишь около 1--10\% от объёма всего образа. Иными словами, в типичном сценарии работы сканера через registry около 90\% скачиваемых данных в конечном итоге оказываются избыточными.

\subsubsection{Влияние на инфраструктуру}

В описанной модели сосредоточен ряд негативных эффектов с точки зрения как самого сканирования, так и инфраструктуры в целом:

\begin{itemize}
\item \textit{сетевой трафик}~--- объем данных, передаваемых между registry и сканером, в значительной части является избыточным, что увеличивает нагрузку на сеть;
\item \textit{нагрузка на сам registry}~--- передача больших блобов требует ресурсов как сети, так и хранилища;
\item \textit{время сканирования}~--- скачивание и распаковка слоев занимают значительную часть времени работы сканера, особенно при сканировании <<с холодного старта>>, когда локальный кеш пуст.
\end{itemize}

Так, при активном сканировании эти эффекты накапливаются и становятся существенным фактором эффективности и стоимости всей инфраструктуры.

\subsection{Формулировка проблемы}

Проведенный обзор позволяет сформулировать основную проблему, рассматриваемую в данной работе:

\textit{Существующие сканеры уязвимостей для анализа образа загружают требуемые слои целиком, в то время как для проведения индексирования требуется лишь малая часть содержимого~--- метаданные пакетных менеджеров и манифесты языковых зависимостей. Это приводит к избыточному сетевому трафику и временным затратам, что особенно критично в сценариях интенсивной эксплуатации.}

Решение этой проблемы естественным образом разделяется на две независимые подзадачи:

\begin{itemize}
\item обеспечить со стороны registry и формата хранения образа возможность отдавать содержимое слоев с гранулярностью, достаточной для эффективного сканирования;
\item поддержать эту возможность со стороны сканера, обеспечив скачивание только необходимыых для анализа фрагментов слоя.
\end{itemize}

\newpage
